\documentclass[12pt]{article}
\usepackage[letterpaper,top=0.75in,bottom=0.75in,left=0.78in,right=1.55in,marginparwidth=1.15in,marginparsep=0.18in]{geometry}
\usepackage{fontspec}
\setmainfont{Latin Modern Roman}
\usepackage{microtype}
\usepackage{fancyhdr}
\usepackage{titlesec}
\usepackage{enumitem}
\usepackage{xcolor}
\usepackage[hidelinks]{hyperref}
\setlength{\parindent}{1.05em}
\setlength{\parskip}{0.32em}
\setlength{\headheight}{15pt}
\titleformat{\section}{\large\bfseries\scshape}{}{0pt}{}
\titleformat{\subsection}{\normalsize\bfseries}{}{0pt}{}
\pagestyle{fancy}
\fancyhf{}
\fancyhead[L]{Whately Logic: Macbeth Lit-Anchor}
\fancyhead[R]{Lesson 2}
\fancyfoot[C]{\thepage}
\renewcommand{\headrulewidth}{0.4pt}

\newcommand{\focusbox}[1]{\par\vspace{0.35em}\noindent\begingroup\setlength{\fboxsep}{6pt}\colorbox{gray!12}{\begin{minipage}{\dimexpr\linewidth-2\fboxsep\relax}#1\end{minipage}}\endgroup\par\vspace{0.45em}}
\newcommand{\studentline}{\par\vspace{0.68em}\hrule\vspace{0.72em}}
\newcommand{\gloss}[1]{\marginpar{\scriptsize\raggedright\setlength{\baselineskip}{7.8pt}\color{gray!70!black}#1}}
\newcommand{\speaker}[1]{\par\vspace{0.35em}\noindent\textsc{#1}\par\vspace{-0.15em}}

\begin{document}
\begin{center}
{\LARGE\bfseries Macbeth Lit-Anchor}\par\vspace{0.18em}
{\large\scshape Words That Lie Like Truth}\par\vspace{0.35em}
{Lesson 2: Ambiguity, Equivocation, and Unstable Terms}\par\vspace{0.55em}
\rule{0.82\textwidth}{0.4pt}
\end{center}

\section*{Purpose}
A lit-anchor connects the week's logic lesson to a recurring literary text. In Lesson 2, Whately teaches that reasoning becomes unstable when a key term can carry more than one meaning. \textit{Macbeth} gives us a terrifying example. The apparitions do not merely tell Macbeth simple lies. They use words that can be heard in one sense while meaning something narrower, stranger, or more dangerous. Macbeth treats unclear language as if it were clear, and his false confidence grows from that mistake.

\focusbox{\textbf{Whately's Lesson 2 habit:} Before testing an argument, fix the meaning of its key terms. If a word or phrase shifts meaning, the conclusion may only appear to follow.}

\section*{Central Question}
When Macbeth trusts words that can mean more than one thing, how does unclear language make false confidence seem reasonable?

\section*{Guided Text: Act 4, Scene 1}
Read the original text below. The side notes are light glosses only; they do not replace the text. Watch especially for phrases that sound certain but are not as clear as Macbeth thinks.

\speaker{First Apparition}
\begin{verse}
Macbeth! Macbeth! Macbeth! beware Macduff;\\
Beware the thane of Fife. Dismiss me. Enough.\gloss{The first warning is direct: beware Macduff. Macbeth should not forget it.}
\end{verse}

\speaker{Macbeth}
\begin{verse}
Whate'er thou art, for thy good caution, thanks;\\
Thou hast harp'd my fear aright: but one word more,--
\end{verse}

\speaker{First Witch}
\begin{verse}
He will not be commanded: here's another,\\
More potent than the first.
\end{verse}

\speaker{Second Apparition}
\begin{verse}
Macbeth! Macbeth! Macbeth!
\end{verse}

\speaker{Macbeth}
\begin{verse}
Had I three ears, I'ld hear thee.
\end{verse}

\speaker{Second Apparition}
\begin{verse}
Be bloody, bold, and resolute; laugh to scorn\\
The power of man, for none of woman born\\
Shall harm Macbeth.\gloss{The phrase sounds universal. Macbeth hears it as if it means ``no human being.''}
\end{verse}

\speaker{Macbeth}
\begin{verse}
Then live, Macduff: what need I fear of thee?\\
But yet I'll make assurance double sure,\\
And take a bond of fate: thou shalt not live;\\
That I may tell pale-hearted fear it lies,\\
And sleep in spite of thunder.\gloss{Macbeth draws a confident conclusion, yet still wants extra safety. His reasoning and fear do not fully agree.}
\end{verse}

\speaker{Third Apparition}
\begin{verse}
Be lion-mettled, proud; and take no care\\
Who chafes, who frets, or where conspirers are:\\
Macbeth shall never vanquish'd be until\\
Great Birnam wood to high Dunsinane hill\\
Shall come against him.\gloss{The words appear impossible if ``wood'' means a whole forest walking by itself.}
\end{verse}

\speaker{Macbeth}
\begin{verse}
That will never be\\
Who can impress the forest, bid the tree\\
Unfix his earth-bound root? Sweet bodements! good!\\
Rebellion's head, rise never till the wood\\
Of Birnam rise, and our high-placed Macbeth\\
Shall live the lease of nature, pay his breath\\
To time and mortal custom.\gloss{Macbeth fixes the meaning too quickly. He decides that the prophecy promises ordinary safety and long life.}
\end{verse}

\section*{Guided Analysis}
The apparitions' words sound like protection. Macbeth turns them into an argument for safety. Whately would ask us to slow the argument down and fix the key terms before trusting the conclusion.

\subsection*{Claim 1: Macbeth is safe from Macduff}
\begin{itemize}[leftmargin=*]
\item \textbf{Macbeth's conclusion:} Macduff cannot harm me.
\item \textbf{Reason:} ``None of woman born / Shall harm Macbeth.''
\item \textbf{Key term or phrase:} ``of woman born.''
\item \textbf{Macbeth's assumed meaning:} every human being; anyone who has a mother.
\item \textbf{Hidden assumption:} the phrase is broad enough to include Macduff and every possible opponent.
\item \textbf{Whately test:} The conclusion depends on an unfixed phrase. Macbeth hears the words in the broadest possible sense, but the phrase may have a narrower meaning.
\end{itemize}

\subsection*{Claim 2: Macbeth cannot be defeated}
\begin{itemize}[leftmargin=*]
\item \textbf{Macbeth's conclusion:} Rebellion cannot defeat me.
\item \textbf{Reason:} Macbeth will not be defeated until Birnam Wood comes to Dunsinane.
\item \textbf{Key term or phrase:} ``Birnam wood ... shall come.''
\item \textbf{Macbeth's assumed meaning:} an actual forest must uproot itself and walk.
\item \textbf{Hidden assumption:} there is no other reasonable way for a wood to ``come'' against a castle.
\item \textbf{Whately test:} Macbeth treats an image as if it has only one literal meaning. The words allow more than his first interpretation.
\end{itemize}

\focusbox{\textbf{Lesson 2 takeaway:} Macbeth is not only deceived by lies. He is deceived by words that can be taken two ways. The argument looks safe because the terms have not been fixed.}

\section*{Discussion}
Use these questions to prepare for the independent passage.

\begin{enumerate}[leftmargin=*]
\item Which warning in the passage is clearest? Which promises are more ambiguous?
\item Why does the phrase ``none of woman born'' sound stronger than it really is?
\item What does Macbeth assume the phrase ``Birnam wood ... shall come'' must mean?
\item How does Macbeth's desire for safety affect the way he interprets the apparitions' words?
\item What would a careful Whately-style reader ask before accepting Macbeth's conclusion?
\end{enumerate}

\section*{Independent Text: Act 5, Scene 8}
Now read the moment when one of the apparitions' phrases is tested. Macbeth finally discovers that the words he trusted were not as stable as he thought.

\speaker{Macbeth}
\begin{verse}
Thou losest labour:\\
As easy mayst thou the intrenchant air\\
With thy keen sword impress as make me bleed:\\
Let fall thy blade on vulnerable crests;\\
I bear a charmed life, which must not yield,\\
To one of woman born.\gloss{Macbeth repeats the phrase as a shield. His confidence depends on his interpretation of ``born.''}
\end{verse}

\speaker{Macduff}
\begin{verse}
Despair thy charm;\\
And let the angel whom thou still hast served\\
Tell thee, Macduff was from his mother's womb\\
Untimely ripp'd.\gloss{Macduff reveals a different sense of ``born'': not delivered in the ordinary manner.}
\end{verse}

\speaker{Macbeth}
\begin{verse}
Accursed be that tongue that tells me so,\\
For it hath cow'd my better part of man!\\
And be these juggling fiends no more believed,\\
That palter with us in a double sense;\\
That keep the word of promise to our ear,\\
And break it to our hope. I'll not fight with thee.\gloss{Macbeth names the trick: a promise heard one way, fulfilled another way.}
\end{verse}

\section*{Independent Analysis}
Answer in complete thoughts. Leave yourself enough evidence that this work can become a portfolio artifact.

\begin{enumerate}[leftmargin=*]
\item What exact phrase from Act 4 does Macbeth repeat in Act 5?\studentline\studentline
\item What meaning did Macbeth give that phrase when he first heard it?\studentline\studentline
\item What alternate meaning does Macduff's revelation expose?\studentline\studentline
\item What hidden assumption made Macbeth's reasoning look stronger than it was?\studentline\studentline
\item Explain Macbeth's line: ``That palter with us in a double sense.'' What does he finally understand about the words he trusted?\studentline\studentline\studentline
\item Use Whately's Lesson 2 habit: how would fixing the meaning of ``born'' have changed Macbeth's confidence?\studentline\studentline\studentline
\end{enumerate}

\section*{Portfolio Artifact: Macbeth Equivocation Chart}
Choose one dangerous phrase from the play. Then complete the chart below.

\begin{itemize}[leftmargin=*]
\item ``of woman born''
\item ``Birnam wood ... shall come''
\item ``fair is foul, and foul is fair''
\item ``safe''
\item ``king''
\item ``chance''
\end{itemize}

\textbf{Phrase:}\studentline
\textbf{Macbeth's interpretation:}\studentline\studentline
\textbf{Another possible meaning:}\studentline\studentline
\textbf{Hidden assumption:}\studentline\studentline
\textbf{How the ambiguity changes the reasoning:}\studentline\studentline\studentline
\textbf{Connection to Whately:}\studentline\studentline

\section*{Closing Synthesis}
Whately teaches that reasoning depends on stable terms. Macbeth shows what happens when a person treats unstable words as if they were clear. The witches and apparitions do not need Macbeth to abandon logic entirely. They only need him to reason from words he has not examined carefully enough. Once the terms shift, Macbeth's conclusion collapses: the word of promise was kept to his ear and broken to his hope.

\end{document}
